% ============================================================
% 06-discussion.tex — V13 Discussion
% ============================================================
\section{Discussion}
\label{sec:discussion}

%------------------------------------------------
\subsection{Architectural Findings and Operational Trade-offs}

\textbf{RQ1 (Architecture-Level Effectiveness):}
The architecture-level comparison indicates that supervisor-routed multi-agent decomposition achieves a 10.0\,pp higher Fact-Grounded Pass Rate ($\text{Pass}_{\ge 50\%}$) point estimate than the Unified Single Agent (76.0\% vs.\ 66.0\%; paired bootstrap CI $[+2.0, +19.0]$\,pp; exact McNemar $p{=}0.0414$), alongside a +3.58\,pp Fact Coverage advantage (95\%\,CI $[-1.2, +8.5]$\,pp) and 41.4\% fewer input tokens (3{,}690 vs.\ 6{,}296), at the expense of higher wall-clock latency (10{,}652 vs.\ 5{,}159\,ms). Routed Generic records lower Fact Coverage than \system{} (54.2\% vs.\ 61.0\%; $\Delta{=}-6.85$\,pp, $p{<}0.05$), indicating that specialist policies contribute substantially to answer fidelity under shared routing.

\textbf{RQ2 \& RQ3 (Mechanisms, Tool Isolation, and Registry Scaling):}
Component ablations reveal two distinct functional roles. \emph{Accuracy drivers:} Heterogeneous evidence access (S2-A3) prevents severe retrieval failure ($-12.3$\,pp Source Recall, $-17.8$\,pp on finance); specialist policies (S2-A1) govern answer coverage ($-6.85$\,pp Fact Coverage); route repair (S2-A5) provides cross-domain robustness ($-15.0$\,pp pass rate); and deterministic calculators (S2-A4) eliminate arithmetic hallucination ($-55.6$\,pp tuition exact match).
Conversely, \emph{scalability and efficiency mechanisms:} At the baseline 11-tool scale, tool isolation (S2-A2) does not affect pass rate ($\Delta\text{Pass}_{.50} = 0.0$\,pp, CI $[-6.0, +6.0]$\,pp) but expands input context by $+1{,}270$ tokens (+34.4\%). This demonstrates that at small registry scales, tool scoping acts as an efficiency and token-containment mechanism rather than an accuracy driver.
Its critical value emerges under registry expansion (Scenario~3): when scaling to 51 tools with semantic distractors, the Routed Specialist sustains 88.0\% tool E2E Success with 9.8 visible tools, outperforming Global Top-10 (78.7\%, $+9.3$\,pp, CI $[+2.0, +17.3]$\,pp; limited by candidate truncation) and Full-Registry Single Agent (83.3\%, $+4.7$\,pp, CI $[-1.3, +11.3]$\,pp) while saving 52.4\% input tokens (1{,}770 vs.\ 3{,}722). Specialization effectively decouples per-agent context load and interference risk from aggregate institutional tool growth.

\textbf{Latency Decomposition and Operational Guidance:}
The empirical data resolve operational trade-offs across architectures:
(1)~\emph{Measurement boundaries:} As detailed in Section~\ref{sec:exp_setup}, the single-agent baseline measures inference post-retrieval ($5{,}159$\,ms), whereas multi-agent pipelines measure end-to-end dispatch inclusive of routing ($8{,}921$--$10{,}652$\,ms).
(2)~\emph{Contract overhead:} The $+1{,}731$\,ms delta between generic and specialist multi-agent execution is primarily driven by structured contract generation (+48\% output tokens, $2{,}092$ vs.\ $1{,}415$), with supervisor routing contributing ${\sim}650$ input tokens.
(3)~\emph{Deployment guidance:} Monolithic single-agent designs remain preferable for low-latency, single-domain workflows with small tool pools (${\le}10$); supervisor-routed multi-agent architectures are warranted when university counseling requires heterogeneous evidence representations, deterministic arithmetic fidelity, and tool isolation against registry expansion.

%------------------------------------------------
\subsection{Relation to Prior Work}

While ReAct~\cite{yao2023react} and Toolformer~\cite{schick2023toolformer} establish tool use within single-agent paradigms, our results indicate that full-registry single agents suffer substantial prompt expansion as tools accumulate. Unlike collaborative multi-agent frameworks (e.g., AutoGen~\cite{wu2023autogen}) where unstructured chatter compounds token consumption and latency, \system{} enforces a bounded, single-owner dispatch model with deterministic route repair. The S2-A3 ablation highlights the fragility of uniform text retrieval on structured relational queries, corroborating the rationale for heterogeneous evidence representations in GraphRAG~\cite{edge2024graphrag}. Furthermore, whereas Adaptive-RAG~\cite{jeong2024adaptiverag} dynamically branches based on linguistic complexity, \system{} routes queries according to institutional data modalities.

%------------------------------------------------
\subsection{Threats to Validity}
\label{sec:limitations}
\textbf{Internal:} All primary evaluations utilize Gemini 2.5 Flash Lite with greedy decoding ($\tau{=}0$) to reduce sampling variance; API serving latency and backend scheduling variations remain potential operational confounds. \textbf{External:} The benchmark is grounded in statutes and curricula from Can Tho University; validation across multi-campus institutions with bilingual or trimester policies remains future work. Semantic-neighbor tools in Scenario~3 are controlled synthetic constructs modeling intra-domain ambiguity. \textbf{Construct:} Metric scoring incorporates automated evaluation via Ragas alongside deterministic exact-match checks; human verification was conducted on a subset to validate metric concordance. \textbf{Statistical:} Evaluations are based on 100 held-out cases with paired bootstrap resampling (10,000 iterations); the arithmetic sub-cohort ($N{=}9$) has a small sample size, though the difference in exact match (8/9 vs.\ 3/9) is substantial. \textbf{Operational:} Latency measurements reflect cloud-hosted API execution and public network connectivity.
